%% ============================================================================
\section{Pipeline Defects Identified During Reproduction}
\label{app:defects}
\addcontentsline{lot}{section}{Pipeline Defects Identified During Reproduction}
%% ============================================================================

Porting the training code to a second environment surfaced a number of defects in
the experimental pipeline.
They are documented here because several are silent - they degrade results or
invalidate measurements without raising an error - and because two of them
invalidate earlier data-staging measurements and therefore bear directly on the
claims this thesis is able to make.

\subsection{Defects affecting the staging mechanism}

\subsubsection{Policy selected after the datasets it selects were constructed}
\label{sec:defect-ordering}

The staging policy was applied only to \code{DataLoader} construction, not to
dataset construction.
In the original control flow, the three datasets were built before the profiler ran:

\begin{lstlisting}[language=Python]
policy = base_config.get("data_policy", "loose")   # key absent -> "loose"
train_dataset = create_dataset(..., policy=policy) # always the loose reader
...
base_config["data_policy"] = policy_summary["policy"]  # profiler decides, too late
train_loader = build_dataloader(train_dataset, policy=base_config["data_policy"])
\end{lstlisting}

The datasets were never rebuilt after the profiler had spoken, and the selected
policy reached only \code{build\_dataloader}, whose entire use of it is to choose
\code{num\_workers} and \code{prefetch\_factor}.
The shard and stage reader classes were therefore unreachable code, and every run
- irrespective of the policy the profiler reported - read loose files.

The consequence is that any previously reported staging result describes a policy
\emph{decision} rather than a policy \emph{execution}.
The fix reorders policy selection ahead of dataset construction, with the profiler
probing an explicitly-loose dataset and the working datasets built from its verdict.

\subsubsection{No archive the shard reader could open}

The shard reader addresses members by original filename within archives named
\code{originals.tar} and \code{negs.tar}.
The repository's shard generator produces WebDataset-format shards
(\code{shard-000000.tar}, members keyed \code{train\_000123.jpg}), which that reader
can neither open nor address; and no archives of either form existed on disk.
The shard branch would therefore have fallen through to its loose fallback even had
it been reachable.

Archives in the required layout are now produced by a dedicated builder, which
verifies its output by comparing decoded samples against the loose path.
They are written uncompressed for the reason given in
Section~\ref{ch:methodology}: the payload is already-compressed JPEG and PNG, so
compression yields negligible size reduction while forcing decompression from the
stream start on every random access.

\subsubsection{Staging silently degraded to a plain read}

The staging dataset copies samples to a scratch directory but swallows
\code{IOError} and \code{OSError} during the copy.
The scratch directory was never created, so every copy failed silently and the
policy read directly from the source path - appearing selected and active while
having no effect.

\subsection{Defects affecting segmentation accuracy}

\begin{table}[!ht]
\centering
\begin{tabular}{p{4.2cm}p{9.2cm}}
\toprule
Defect & Description and evidence \\
\midrule
Focal loss $\alpha$ unset in one variant
  & \code{smp.losses.FocalLoss} defaults to \code{alpha=None}, applying no
    class weighting on a target with a median 7\% positive-pixel rate. \\
Validation on a random crop
  & Scoring a different random patch of each tile every epoch injects variance into
    checkpoint selection and early stopping. Note that the reference
    \emph{deliberately} random-crops at test time and averages ten passes; the
    defect is using it as an unaveraged per-epoch selection signal. \\
Configured loss silently ignored
  & One model hardcoded \code{BCEWithLogitsLoss} while both the configuration and
    the training log reported \code{focal}. Correcting it moved epoch-1 IoU from
    0.299 to 0.362 and reversed a declining validation trend. Unknown loss names
    now raise rather than falling through. \\
fp16 divergence on EfficientNet
  & Under \code{16-mixed}, training loss became non-finite at epoch 4 and all
    metrics collapsed to zero, while early stopping continued for its full
    15-epoch patience on a dead model; the run reported a ``best'' score that was
    merely the last value before the blowup. ResNet-50 was unaffected across 20
    epochs. Isolating the loss under \code{autocast} showed gradients bit-identical
    to fp32, locating the overflow in the encoder - consistent with
    EfficientNet's SiLU and squeeze-excite blocks exceeding fp16's exponent range.
    Resolved by \code{bf16-mixed}; a guard now aborts at the epoch a loss becomes
    non-finite. \\
IoU averaging convention (main pipeline)
  \label{sec:defect-metric}
  & \code{PADS/*/model.py} computes validation/test IoU with
    \code{torchmetrics.BinaryJaccardIndex}, which pools confusion-matrix counts
    across the whole split before dividing, rather than the reference's
    per-image mean (Section~\ref{sec:repro-metric}, which quantifies a 3.7-point
    effect from this convention alone on one checkpoint). Not present in either
    the Appendix~\ref{app:reproduction} workstation port or
    \code{casini\_reference}, both of which use \code{smp.metrics.iou\_score}
    with \code{reduction="macro-imagewise"} throughout. \\
Flip-augmentation semantics (main pipeline)
  \label{sec:defect-augmentation}
  & \code{PADS/*/dataset.py} applies two independent single-axis flip
    transforms rather than the reference's \code{A.Flip}, which selects one of
    \{vertical, horizontal, both\} uniformly (Appendix~\ref{app:reproduction},
    Section~\ref{sec:repro-faithfulness} measures the reference semantics
    empirically at 0.752 no-flip and 0.089/0.084/0.075 for the three flip modes,
    against an expected 0.750/0.083 each). Changes the effective
    augmentation distribution seen during training; not quantified in isolation
    against IoU, since it is confounded with the metric-convention defect above
    in every main-pipeline run collected this session. \\
\bottomrule
\end{tabular}
\caption{Accuracy-affecting defects. The final two rows are specific to
\code{PADS/*/model.py} and \code{PADS/*/dataset.py} (Sections~\ref{ch:results}
and~\ref{ch:ablation}'s pipeline); neither is present in
Appendix~\ref{app:reproduction}'s workstation port or in
\code{casini\_reference} (Section~\ref{sec:repro-closure}).}
\label{tab:accuracy-defects}
\end{table}

\subsection{Defects affecting reproducibility and operations}

\begin{itemize}
  \item \textbf{Metric-name path separator.} Checkpoint filenames interpolate the
    monitored metric name \code{valid/iou} verbatim. On filesystems treating
    \code{/} as a path separator this creates a directory rather than a filename,
    and \code{save\_top\_k} then prunes the checkpoint while leaving an orphaned
    directory per improved epoch.
  \item \textbf{Profiler perturbs the training RNG.} The warm-up pass iterates the
    \emph{training} dataset, whose transforms are stochastic, advancing the global
    random state before training begins. A profiled run is therefore not
    seed-comparable with an unprofiled one at identical configuration. Any
    comparison between them must either restore RNG state after profiling or treat
    the difference as within run-to-run variance, which was measured at
    approximately $\pm 0.02$ IoU.
  \item \textbf{Study persistence.} Hyperparameter studies were not persisted, so an
    interrupted sweep restarted from trial zero and reused checkpoint directories.
\end{itemize}

\subsection{Defects identified while collecting the Section~\ref{ch:results} and~\ref{ch:ablation} evidence}
\label{sec:defects-session2}

Four further defects were identified while running the full evaluation sweep.
The first two are infrastructure failures rather than correctness bugs, but are
documented because they caused real job losses and shaped the retry pattern
visible in some of the job identifiers referenced in Sections~\ref{ch:results}
and~\ref{ch:ablation}.

\begin{itemize}
  \item \textbf{Scratch directory scoped to \code{project\_root}, not shared.}
    \code{resolve\_scratch\_dir()} derived the node-local staging path from each
    job's \code{project\_root}, which is deliberately unique per experiment
    configuration for checkpoint isolation. Because every \code{stage}-policy
    job across the evaluation sweep therefore had its own scratch directory, each
    of the roughly forty such jobs independently copied its own multi-gigabyte
    snapshot of the dataset, accumulating 69 redundant copies ($\approx$91\,GB)
    and exhausting the account's disk quota mid-sweep, which in turn failed every
    subsequent job attempting to write output. A real node-local scratch mount is
    shared per node regardless of which job is using it; the fix decouples
    \code{scratch\_dir} from \code{project\_root} and defaults to one fixed,
    shared location unless explicitly overridden.
  \item \textbf{Checkpoint auto-resume from a truncated file.} Several jobs that
    had progressed far enough to begin writing \code{last.ckpt} before the disk
    quota was exhausted left a truncated checkpoint behind. Lightning resumes
    training from \code{last.ckpt} by default if one is present in the
    checkpoint directory, and loading a truncated \code{.ckpt} raises
    \code{RuntimeError: PytorchStreamReader failed reading zip archive} rather
    than falling back to a fresh run. Affected job directories were identified
    and their partial checkpoints removed before resubmission.
  \item \textbf{\code{casini\_reference}'s epoch-end hooks never wired in for
    two of three architectures.} \code{build\_model()} and
    \code{load\_model\_from\_checkpoint()} instantiated the plain
    \code{ArcheoModel} base class for the MA-Net and DeepLabV3+ code paths,
    rather than \code{\_EpochBufferedArcheoModel}, the subclass that implements
    \code{on\_\{train,validation,test\}\_epoch\_end} in place of the
    Lightning~1.x \code{*\_epoch\_end(self, outputs)} hooks the ported notebook
    code assumes (Lightning~2.x removed these; see also
    Section~\ref{sec:repro-faithfulness}). Only the SegFormer path used the
    correct subclass. Training completed without error under the plain base
    class -- per-step loss logging does not depend on the missing hook -- so the
    defect was silent until the test loop's aggregated metric lookup raised
    \code{KeyError: 'test/IOU-img'}, since the epoch-aggregated IoU had never
    actually been logged. Both non-SegFormer paths now use the buffered
    subclass; this was confirmed by a full run reaching the expected metric
    before the Table~\ref{tab:repro-full} results were collected.
  \item \textbf{Profiler warm-up window bias -- implementation diverged from its
    own documented design.} Documented in full in
    Section~\ref{sec:res-profiler-bias} rather than repeated here: the default
    8-batch profiling window included one-time CUDA/cuDNN warm-up cost in its
    \Tgpu{} average, inflating it 6--10$\times$ and systematically suppressing
    the adaptive policy across most of Section~\ref{ch:ablation}'s threshold
    sweep. This is not merely an undocumented gap: Section~\ref{ch:methodology}
    already specifies discarding the first $w_0=10$ batches and aggregating by
    median rather than mean, for exactly this reason, in exactly these words.
    The shipped \code{profile\_data\_pipeline()} implemented neither. Discard
    was implemented and validated this session (\code{--sweep\_discard}); median
    aggregation was not, and remains genuinely open -- it is plausible that
    switching from mean to median would provide some robustness independent of
    discarding, since a median is inherently less sensitive to a small number
    of outlying (warm-up) samples than a mean is, but this has not been tested,
    and claiming it as fixed would overstate what this session actually did.
    Included in this list because, unlike the first three defects, it is a
    correctness defect in \pads{} itself rather than in the evaluation
    infrastructure, and it directly affects how several results in
    Sections~\ref{ch:results} and~\ref{ch:ablation} should be interpreted.
\end{itemize}

\subsection{Implication for prior results}

Sections~\ref{sec:defect-ordering} and the missing archives are the consequential
ones.
Because the selected policy never reached the data path, and because the shard
policy had no readable data in any case, staging results collected before these
fixes report decisions rather than executions.
They are not recoverable by reanalysis and the corresponding experiments require
repetition.
The defects are properties of the code rather than of any particular machine, and
therefore applied equally to cluster runs.
